\documentclass[../../main.tex]{subfiles}

\begin{document}

\section{Practice Problems}

Below are some practice problems with a very subjective rating
one difficulties. The solutions are included at the end of the
Math Olypmiad part of LibreNotebook.

\begin{listproblem}[2004 USAMO Problem 5 \easy]
Show that the following holds for $a, b, c > 0$.
\[
(a^5 - a^2 + 3)(b^5 - b^2 + 3)(c^5 - c^2 + 3)
\geq (a + b + c)^3
\]
\end{listproblem}

\begin{listproblem}[2001 IMO Shortlist A3 \easy]
Show that the following inequality holds for arbitrary reals
$x_1, x_2, \cdots, x_n$.
\[
\frac{x_1}{1 + x_1^2} + \frac{x_2}{1 + x_1^2 + x_2^2} +
    \cdots + \frac{x_n}{1+ x_1^2  + \cdots + x_n^2}
< \sqrt{n}
\]
\end{listproblem}

\begin{listproblem}[2025 KMO Problem 6 \medium]
For arbitrary reals $x_1, x_2, \dots, x_{99}$ that satisfy
$0 < x_1 < x_2 < \cdots < x_{99}$, determine the minimum value of
$c > 0$ such that the following inequality always holds.
\begin{align*}
    &3\sqrt{x_1} + 4 \left( \sqrt{x_2 - x_1} + \sqrt{x_3 - x_2} +
        \cdots + \sqrt{x_{99} - x_{98}}
    \right) \\
    &\qquad\qquad\qquad\qquad\qquad
    \leq c \left(
        \sqrt{x_2} + \sqrt{x_3} + \cdots + \sqrt{x_{98}}
    \right) + 5\sqrt{x_{99}}
\end{align*}
\end{listproblem}

\begin{listproblem}[2023 IMO Problem 4 \medium]
Let $x_1, x_2, \cdots , x_{2023} > 0$ be pairwise different numbers
such that $a_n$ defined as the following is an integer for all
$n = 1, 2, \cdots, 2023$.
\[
a_n
= \sqrt{ \left( x_1 + x_2 + \cdots + x_n \right) \left( \frac{1}{x_1}
    + \frac{1}{x_2} + \cdots + \frac{1}{x_n} \right)
}
\] 
Show that $a_{2023} \geq 3034$.
\end{listproblem}






\end{document}


